\documentclass{article}%
\usepackage[T1]{fontenc}%
\usepackage[utf8]{inputenc}%
\usepackage{lmodern}%
\usepackage{textcomp}%
\usepackage{lastpage}%
\usepackage{geometry}%
\geometry{margin=2.2cm}%
%
\usepackage{amsmath}%
\usepackage{amssymb}%
\usepackage{booktabs}%
\usepackage{longtable}%
\usepackage{array}%
\usepackage{hyperref}%
\hypersetup{hidelinks}%
\title{IME 2027 {-} Física (sem judge) x Poliedro\\[4pt]\large Comparação com o gabarito Poliedro\thanks{Gerado com gpt-resolve: \url{https://github.com/lgabs/gpt-resolve}}}%
\date{30/10/2026}%
%
\begin{document}%
\normalsize%
\maketitle%
\section*{Resumo}\textbf{Nota média: 10,0/10} (30 de 30 itens com a resolução correta; cada item vale o mesmo dentro da questão).\par\medskip%
\begin{longtable}{@{}c c p{0.78\textwidth}@{}}\toprule Questão & Nota & Resumo \\ \midrule
1 & 10,0 & Questão 1 inteiramente correta e idêntica ao gabarito Poliedro nos três itens: a) 250 m/s, b) 4000 K, c) 864 kJ. Conferi a transcrição com o PDF original e refiz as contas em Python. A resolução vai além do gabarito: confirma que o projeto é coerente (vértice em (2000; 1000) m), verifica que o foguete ainda sobe na vertical do alvo (v\_y $\approx$ 63,6 m/s) e confere o trabalho por um segundo caminho (W = m(c\_p - c\_v)$\Delta$T). Não há divergências; só restam pontos menores de convenção de sinal e de hipóteses de modelagem. \\
2 & 10,0 & A resolução coincide com o gabarito Poliedro nos três itens. No item a, A = 200 kN para cima e B = 100 kN para baixo, sem componentes horizontais. No item b, CD = DE = 75 kN (compressão), CF = EF = 125 kN (tração), DF = AD = 200 kN (compressão), BE = 100 kN (tração) e AE = 0. No item c, AE é a diagonal de contraventamento: tem esforço nulo neste carregamento, mas impede o mecanismo de paralelogramo no painel ADEB. Refiz o equilíbrio nó a nó e resolvi numericamente o sistema 12$\times$12, que tem posto 12. Os valores conferem, e a matriz sem AE tem posto 11, o que confirma que a estrutura seria hipostática. Não há divergência entre resolução e gabarito. \\
3 & 10,0 & Nos três itens, a resolução coincide com o gabarito Poliedro. Em a, a reação é de 300 N, vertical para cima, e a tração inicial também é de 300 N. Em b, $\Delta$$\theta$ = 100 $^\circ$C. Em c, Q = 5$\cdot$10$^{6}$ J e o rompimento ocorre às 37,5\% do dia. Refiz todas as contas: massa de 50 kg, peso de 600 N, condição de momentos com o CM em L'/2 e energia E = 0,05$\tau$$^2$ na rampa. Estão corretas. As hipóteses implícitas (emissão térmica desprezada, temperatura uniforme, efeito desprezível da dilatação da área) são razoáveis e foram declaradas. A estimativa de $\approx$ 5 s para o efeito da dilatação da área foi confirmada numericamente. Não há erro em nenhum dos dois. \\
4 & 10,0 & A resolução está correta e coincide com o gabarito Poliedro nos três itens. Em (a), V\_AB = 900 V, i1 = 50 A e i2 = 0. Em (b), I/I\_F = 2,5 e o fusível rompe em $\approx$1,5 s. Em (c), I = 500 A, razão 25 e t = 100 ms > 50 ms, então o fusível não protege. Conferi as contas em python3 e a leitura do gráfico na figura original. A única ressalva é a incerteza de leitura gráfica no item (b), com t entre $\approx$1,45 e 1,6 s, que não altera a conclusão. \\
5 & 10,0 & A resolução está correta e bate com o gabarito Poliedro nos três itens: a) T = 450 K; b) Q = 90 kW; c) a proposta de B não é realista. Com a água aquecida, B entrega 71,7 kW e não 116,7 kW, o que dá no máximo 215 residências, e $\Delta$$\dot S$ $\approx$ +0,04 kW/K. No item c, a resolução ainda mostra que o total pedido (195 kW) viola a 1ª lei e que a 2ª lei limita $\dot W$ a 85 kW. O único problema é um erro de digitação na desigualdade de Clausius do item c: o sinal está invertido em relação à convenção que a própria resolução declarou, embora a conclusão (Q$_{3}$ $\geq$ 20 kW) esteja certa. Isso pode render um desconto pequeno de forma, mas não muda nenhum resultado. \\
6 & 10,0 & A resolução está certa e bate com o gabarito Poliedro nos quatro itens. Os resultados são v = 68 m/s, R = 5$\surd$3 m $\approx$ 8,66 m, T $\approx$ 1,07$\cdot$10$^3$ N e V\_sub $\approx$ 5,51$\cdot$10$^{-}$$^2$ m$^3$ $\approx$ 55$\cdot$10$^{-}$$^3$ m$^3$. As duas usam o mesmo caminho: Doppler em A com a velocidade apontando direto para O, e em B com a projeção v$\cdot$cos$\theta$ = 34 m/s, que dá $\theta$ = 60$^\circ$ e R = ($\surd$3/2)$\cdot$d. Depois vêm T$\cdot$sen$\alpha$ = mv$^2$/R e o equilíbrio vertical E = mg + T$\cdot$cos$\alpha$. Conferi o enunciado no PDF original e refiz as contas em Python: tudo bate. As diferenças são só de arredondamento e notação. Não encontrei erro em nenhum dos dois. \\
7 & 10,0 & A resolução e o gabarito Poliedro dão as mesmas respostas nos três itens: a) d = 0,45 m; b) razão $\approx$ -1,6, com kx $\approx$ 6,44 N; c) v = $\surd$2 m/s quando P2 chega ao vértice C. Refiz as contas em python e todas conferem. O ponto discutível é o item b: os dois supõem equilíbrio quase estático de P2, e com a mola partindo relaxada e sem amortecimento isso é uma idealização. O enunciado pede só uma estimativa, então essa hipótese é aceitável. A resolução está correta e completa. \\
8 & 10,0 & A resolução está correta e coincide com o gabarito Poliedro nos quatro itens. Em (a), o anteparo se afasta. Em (b), D(t1) = 0,2$\surd$7 $\approx$ 0,53 m. Em (c), v $\approx$ 0,254 m/s. Em (d), $\Delta$s $\approx$ 1,79 m (1,792 m). A transcrição confere com o PDF original. Refiz as contas em Python com as raízes da tabela e com raízes exatas, e as únicas diferenças são de arredondamento. \\
9 & 10,0 & A resolução está correta e coincide com o gabarito Poliedro nos dois itens: a) u = c($\surd$3h - L)/($\surd$3L - h); b) L/$\surd$3 < h < L. Refiz o sistema de impacto e conferi numericamente: com L = 1 e h = 0,8, o raio refletido sai exatamente na direção -y. A resolução vai além do gabarito ao provar que o impacto encontrado é o primeiro contato e ao mostrar que todo o intervalo do item b é atingível. Não encontrei erro em nenhum dos dois. \\
10 & 10,0 & A resolução e o gabarito do Poliedro concordam nos dois itens. Em (a), só a situação 2 serve, com as mesmas justificativas para descartar 1, 3 e 4. Em (b), ambos chegam a p = ($\pi$d/2)$\surd$($\pi$mv/($\mu$0 q i)), e o período T = $\pi$d$\surd$($\pi$m/($\mu$0 q i v)) é o mesmo. Refiz a física (sentidos dos campos, sinal de F\_x, constante K = 4$\mu$0 q i v/($\pi$d$^2$)) e conferi com uma simulação numérica da trajetória exata: o passo simulado bate com a fórmula até a 5ª casa. A instabilidade da situação 3 também se confirmou na simulação. Não encontrei erro: os dois estão corretos. \\
\bottomrule\end{longtable}%
\section*{Principais pontos de dúvida}%
\begin{itemize}%
\item%
\textbf{Q5} \textit{(solution)} No item c, a desigualdade de Clausius está escrita com o sinal trocado. As hipóteses adotam $\Sigma$ Q\_i/T\_i $\leq$ 0, com Q\_i positivo quando o calor entra no ciclo; isso daria +150/900 - 45/450 - Q$_{3}$/300 $\leq$ 0. A resolução, porém, escreve -150/900 + 45/450 + Q$_{3}$/300 $\leq$ 0, que é a variação de entropia dos reservatórios (deveria ser $\geq$ 0). Do jeito que está escrita, a desigualdade levaria a Q$_{3}$ $\leq$ 20 kW, mas a resolução conclui Q$_{3}$ $\geq$ 20 kW, que é o correto. É um erro de digitação que inverte o sentido da desigualdade. Não afeta a resposta final ($\dot W$ $\leq$ 85 kW), mas um corretor rigoroso pode descontar um pouco. A mesma mistura de convenções aparece na verificação do item b, onde se usa -150/900 + 90/450 > 0 como geração de entropia.%
\end{itemize}%
\section*{Questão 1}%
\href{https://poliedroresolve.sistemapoliedro.com.br/vestibulares/ime/2027/ime-2-fase-dia-2-27-10-2026-2-fase-dia-2/dissertativa-1-fisica-2-fase-dia-2-ime-2027}{Resolução do gabarito}\par\smallskip%
\begin{longtable}{@{}p{0.06\textwidth} p{0.33\textwidth} p{0.33\textwidth} p{0.1\textwidth} p{0.08\textwidth}@{}}\toprule Item & Resolução & Gabarito & Veredito & Certo \\ \midrule
a & v = 250 m/s (a 45$^\circ$) & v0 = 250 m/s & igual & ambos \\
b & T = 4000 K & T\_falha = 4000 K & igual & ambos \\
c & W = 864 kJ (realizado pelos gases) & W = 864 kJ (Q = 2880 kJ, $\Delta$U = 2016 kJ) & igual & ambos \\
\bottomrule\end{longtable}%
\begin{itemize}%
\item%
\textbf{a)} Pela equação da trajetória a 45$^\circ$, y = x - g x$^2$/v$^2$. Com x = 2000 m e y = 1360 m, sai v$^2$ = 10$\cdot$4$\cdot$10$^{6}$/640 = 62500, logo v = 250 m/s (conferido em Python). A resolução também verifica que o foguete ainda sobe nesse ponto: v\_y = 45$\surd$2 $\approx$ 63,6 m/s > 0, e o vértice fica em x = 3125 m > 2000 m. Isso confirma o enunciado.%
\item%
\textbf{b)} Como T $\propto$ v$^2$, T = 2560$\cdot$(250/200)$^2$ = 2560$\cdot$1,5625 = 4000 K. Logo $\Delta$T = 1440 K.%
\item%
\textbf{c)} Q = m c\_p $\Delta$T = 2,0$\cdot$1,0$\cdot$1440 = 2880 kJ. Pela 1ª lei, W = Q - $\Delta$U = 2880 - 2016 = 864 kJ. A resolução confere o valor de outro jeito: c\_v = 0,70 kJ/(kg$\cdot$K), $\gamma$ = 10/7 $\approx$ 1,43, e W = m(c\_p - c\_v)$\Delta$T = 864 kJ.%
\end{itemize}%
\paragraph{Pontos de dúvida}%
\begin{itemize}%
\item%
\textit{(statement)} O enunciado não diz se o 'trabalho realizado' é pelos gases ou sobre eles. Na expansão isobárica com aquecimento, o natural é o trabalho pelos gases, positivo (+864 kJ). A resolução declara essa convenção e o gabarito usa a mesma ($\Delta$U = Q - W). Uma resposta com o sinal oposto precisaria explicitar a convenção.%
\item%
\textit{(statement)} A relação T $\propto$ v$^2$ é uma hipótese de modelagem imposta pelo enunciado, não algo que se deduz da física do foguete. As duas resoluções a usam como dada. A resolução também supõe, sem dizer, que o foguete vira projétil logo ao sair do lançador. Isso é coerente com 'despreze a resistência do ar' e com o projeto (vértice no alvo), que a resolução verifica.%
\end{itemize}%
\section*{Questão 2}%
\href{https://poliedroresolve.sistemapoliedro.com.br/vestibulares/ime/2027/ime-2-fase-dia-2-27-10-2026-2-fase-dia-2/dissertativa-2-fisica-2-fase-dia-2-ime-2027}{Resolução do gabarito}\par\smallskip%
\begin{longtable}{@{}p{0.06\textwidth} p{0.33\textwidth} p{0.33\textwidth} p{0.1\textwidth} p{0.08\textwidth}@{}}\toprule Item & Resolução & Gabarito & Veredito & Certo \\ \midrule
a & R\_A = 200 kN, vertical, para cima; R\_B = 100 kN, vertical, para baixo (o apoio puxa/ancora a estrutura); A\_x = B\_x = 0 & A = 200 kN vertical para cima; B = 100 kN vertical para baixo (A\_x = B\_x = 0) & igual & ambos \\
b & CD = DE = 75 kN (compressão); CF = EF = 125 kN (tração); DF = AD = 200 kN (compressão); BE = 100 kN (tração); AE = 0 & N\_CD = 75 kN (C); N\_CF = 125 kN (T); N\_DF = 200 kN (C); N\_AD = 200 kN (C); N\_AE = 0; N\_DE = 75 kN (C); N\_BE = 100 kN (T); N\_EF = 125 kN (T) & igual & ambos \\
c & AE é a diagonal de contraventamento: triangula o painel ADEB e impede que ele se comporte como um paralelogramo articulado (mecanismo de deslocamento lateral). Sem ela, b + r = 11 < 12 e a estrutura é hipostática, com equilíbrio instável. Sob esta carga vertical, N\_AE = 0, mas a barra garante a rigidez e resiste a cargas horizontais ou a outras distribuições de carga. & N\_AE = 0 só para este carregamento; AE é necessária para outras distribuições de carga e garante a rigidez (sem ela, ADEB vira um paralelogramo articulado e o guindaste balança lateralmente) & igual & ambos \\
\bottomrule\end{longtable}%
\begin{itemize}%
\item%
\textbf{a)} Refiz a conta. BE é a única barra em B e é vertical, então B\_x = 0 e daí A\_x = 0. Momentos em A: 3$\cdot$P + 3$\cdot$B\_y = 0, logo B\_y = -100 kN. Somando as forças verticais: A\_y = 200 kN. O gabarito chega ao mesmo resultado tomando momentos em B (6P = 3A\_y). A solução numérica do sistema completo de 12 equações dá A\_y = 200, B\_y = -100 e componentes horizontais nulas.%
\item%
\textbf{b)} Os oito esforços coincidem em módulo e no tipo (tração ou compressão). Conferi pelo método dos nós e pela solução numérica da matriz 12$\times$12, que tem posto 12 (estrutura isostática). Ela dá CD = DE = -75, CF = EF = +125, DF = AD = -200, AE = 0 e BE = +100 kN, na convenção tração > 0. As equações de cada nó na resolução estão corretas, inclusive a do nó E, onde 75 - 75 - (3/5)N\_AE = 0 leva a N\_AE = 0. A verificação nos apoios fecha com o item a.%
\item%
\textbf{c)} A conclusão é a mesma: esforço nulo neste carregamento, função de travamento/rigidez e o painel ADEB virando paralelogramo sem a barra. A resolução vai além e acrescenta a contagem b + r = 11 < 12. Conferi numericamente que, sem AE, a matriz fica 12$\times$11 com posto 11, ou seja, há um mecanismo. A resolução também afirma que o equilíbrio seria instável porque a carga desce no deslocamento lateral. Isso está certo: a parte superior translada com y\_C = 4cos$\theta$ + cte, que é máximo em $\theta$ = 0, então a energia potencial da carga é máxima na posição de equilíbrio.%
\end{itemize}%
\paragraph{Pontos de dúvida}%
\begin{itemize}%
\item%
\textit{(statement)} A reação em B é para baixo (100 kN), então o apoio B precisa ser ancorado e resistir a arrancamento. A figura mostra um apoio articulado comum, sem indicar essa capacidade. Tanto a resolução quanto o gabarito supõem, corretamente, que o apoio resiste à tração. A resolução declara isso explicitamente.%
\item%
\textit{(solution)} No item c, a resolução fala em instabilidade (a carga desce no mecanismo), um argumento que o gabarito não usa. O argumento está correto, mas é uma justificativa extra. Não muda a conclusão.%
\end{itemize}%
\section*{Questão 3}%
\href{https://poliedroresolve.sistemapoliedro.com.br/vestibulares/ime/2027/ime-2-fase-dia-2-27-10-2026-2-fase-dia-2/dissertativa-3-fisica-2-fase-dia-2-ime-2027}{Resolução do gabarito}\par\smallskip%
\begin{longtable}{@{}p{0.06\textwidth} p{0.33\textwidth} p{0.33\textwidth} p{0.1\textwidth} p{0.08\textwidth}@{}}\toprule Item & Resolução & Gabarito & Veredito & Certo \\ \midrule
a & N = 300 N, vertical para cima (m = 50 kg, P = 600 N, T0 = 300 N) & T = 300 N; R\_A = 300 N vertical para cima & igual & ambos \\
b & $\Delta$$\theta$ = 100 $^\circ$C (T = T0(1+$\alpha$$\Delta$$\theta$), rompe quando $\alpha$$\Delta$$\theta$ = 10$^{-}$$^3$) & $\Delta$T = 100 $^\circ$C & igual & ambos \\
c & Q = 5$\cdot$10$^{6}$ J; $\tau$ = 10$^{4}$ s = 12,5\% do dia após o amanhecer; rompimento às 37,5\% do dia & Q = 5$\cdot$10$^{6}$ J; rompe às 37,5\% do dia (12,5\% após 25\%) & igual & ambos \\
\bottomrule\end{longtable}%
\begin{itemize}%
\item%
\textbf{a)} V = 2$\cdot$0,25$\cdot$0,05 = 0,025 m$^3$ e m = 50 kg, logo P = 600 N com g = GM/R$^2$ = 12 m/s$^2$. Tomando momentos no apoio: T$\cdot$L = P$\cdot$L/2, o que dá T = 300 N e R = 300 N, vertical para cima. Conferi em python. O valor, o sentido e a justificativa (nenhuma força horizontal, então a reação é vertical) batem com o gabarito.%
\item%
\textbf{b)} O apoio esquerdo fica fixo e a haste continua em x = L. O CM vai para L'/2, com L' = L(1+$\alpha$$\Delta$$\theta$). Daí T$\cdot$L = P$\cdot$L'/2, ou seja, T = 300(1+$\alpha$$\Delta$$\theta$). Igualando a 300,3 N, vem $\alpha$$\Delta$$\theta$ = 10$^{-}$$^3$ e $\Delta$$\theta$ = 100 $^\circ$C. O raciocínio é o mesmo do gabarito.%
\item%
\textbf{c)} Q = 50$\cdot$1000$\cdot$100 = 5$\cdot$10$^{6}$ J e A = 0,5 m$^2$. Na rampa, I = 0,2$\tau$ W/m$^2$, então E = 0,05$\tau$$^2$ = 5$\cdot$10$^{6}$ J e $\tau$ = 10$^{4}$ s, que é 12,5\% de 8$\cdot$10$^{4}$ s. O horário é 37,5\%. A conta do gabarito, com t em \% do dia, é equivalente: 32000t$^2$ = 5$\cdot$10$^{6}$ $\Rightarrow$ t = 12,5. Também simulei numericamente a correção pela dilatação da área (A $\propto$ (1+$\alpha$$\Delta$$\theta$)$^2$): ela antecipa o rompimento em $\approx$ 5,0 s (0,006\% do dia), exatamente o valor que a resolução estimou. Os 0,05\% citados pelo gabarito correspondem à variação relativa do tempo (5 s em 10$^{4}$ s), o que é coerente.%
\end{itemize}%
\paragraph{Pontos de dúvida}%
\begin{itemize}%
\item%
\textit{(solution)} A resolução e o gabarito desprezam a emissão térmica (radiação/convecção) de um objeto 'escuro' que é aquecido em 100 $^\circ$C. Essa hipótese é necessária para chegar ao resultado esperado e a resolução a declara explicitamente. O enunciado não dá emissividade nem temperatura ambiente, então a hipótese é razoável.%
\item%
\textit{(statement)} 'Horário em \% do dia' pode ser lido como o horário absoluto (37,5\%) ou como o intervalo depois do amanhecer (12,5\%). A resolução e o gabarito informam os dois valores. A leitura natural é 37,5\%.%
\item%
\textit{(key)} O gabarito diz que a dilatação superficial 'mudaria só 0,05\%'. A área chega a variar até $\approx$ 0,2\%, mas o tempo de rompimento varia $\approx$ 0,05\% (cerca de 5 s). Isso é só uma questão de redação: as duas estimativas são coerentes entre si e com a simulação numérica.%
\end{itemize}%
\section*{Questão 4}%
\href{https://poliedroresolve.sistemapoliedro.com.br/vestibulares/ime/2027/ime-2-fase-dia-2-27-10-2026-2-fase-dia-2/dissertativa-4-fisica-2-fase-dia-2-ime-2027}{Resolução do gabarito}\par\smallskip%
\begin{longtable}{@{}p{0.06\textwidth} p{0.33\textwidth} p{0.33\textwidth} p{0.1\textwidth} p{0.08\textwidth}@{}}\toprule Item & Resolução & Gabarito & Veredito & Certo \\ \midrule
a & V\_AB = 900 V; i\_Rec1 = 50 A, i\_Rec2 = 0 (I no gerador/fusível = 50 A) & U\_AB = 900 V; i1 = 50 A; i2 = 0 & igual & ambos \\
b & Sim, rompe: I/I\_F = 2,5 > \textasciitilde{}1,1 e t $\approx$ 1,5 s (leitura gráfica) & I/I\_fus = 2,5 $\Rightarrow$ rompe em t $\approx$ 1,5 s & igual & ambos \\
c & I\_cc = 500 A, I/I\_F = 25, t = 0,1 s = 100 ms > 50 ms; o fusível não protege o gerador & I\_curto = 500 A, razão 25 $\Rightarrow$ t = 100 ms > 50 ms $\Rightarrow$ NÃO protege & igual & ambos \\
\bottomrule\end{longtable}%
\begin{itemize}%
\item%
\textbf{a)} Refiz a conta. Pelo teorema de Millman, V = (500+850+450)/2 = 900 V, logo I = (1000-900)/2 = 50 A, i1 = 900-850 = 50 A e i2 = (900-900)/2 = 0. O balanço de potência fecha: 50 kW fornecidos = 5 kW dissipados no gerador + 45 kW absorvidos por Rec1. A resolução e o gabarito chegam aos mesmos valores.%
\item%
\textbf{b)} Razão: 50/20 = 2,5. Medindo a figura original, o joelho da curva fica em cerca de (1,15--1,2; 3,8 s) e o trecho reto termina em cerca de (3,3--3,4; 0,1 s). Testei a interpolação em python3 para todas essas combinações e obtive t entre 1,45 e 1,61 s. Isso é coerente com $\approx$1,5 s dos dois. A resolução também justifica que haverá rompimento, porque 2,5 está à direita da assíntota em $\approx$1,1, como o item pede ('se e quando').%
\item%
\textbf{c)} Com o curto em A--B, V\_AB = 0. As correntes que os receptores injetariam circulam pelo curto e não passam por F, então I\_F = 1000/2 = 500 A e a razão é 25. A curva se mantém no patamar mínimo de 10$^{-}$$^{1}$ s = 100 ms, que é mais que os 50 ms suportados. Logo o fusível não protege. A resolução e o gabarito tiram a mesma conclusão.%
\end{itemize}%
\paragraph{Pontos de dúvida}%
\begin{itemize}%
\item%
\textit{(statement)} Os tempos dos itens (b) e (c) vêm de leitura gráfica. Na figura original, a linha tracejada fica em $\approx$1,1--1,15, e o início do trecho decrescente ($\approx$3,8 s) parece estar em $\approx$1,2, não em 1,1 como diz a transcrição. Dependendo da leitura, t no item (b) fica entre $\approx$1,45 e $\approx$1,6 s. A conclusão não muda, mas o corretor deve aceitar uma faixa em torno de 1,5 s.%
\item%
\textit{(statement)} No item (c), a razão 25 está fora da escala do gráfico, que vai só até 5. Tanto a resolução quanto o gabarito supõem que o patamar de 0,1 s continua para razões maiores. A hipótese é razoável, já que o patamar representa o tempo mínimo de atuação, mas não está explícita no enunciado.%
\item%
\textit{(statement)} No item (a), 'imediatamente após o fechamento' é tratado com as f.c.e.m. dadas, valendo desde t = 0. Num motor real que parte do repouso, a f.c.e.m. inicial seria nula. O enunciado, porém, fornece valores fixos de f.c.e.m., e a resolução e o gabarito adotam esse modelo, que é o esperado.%
\end{itemize}%
\section*{Questão 5}%
\href{https://poliedroresolve.sistemapoliedro.com.br/vestibulares/ime/2027/ime-2-fase-dia-2-27-10-2026-2-fase-dia-2/dissertativa-5-fisica-2-fase-dia-2-ime-2027}{Resolução do gabarito}\par\smallskip%
\begin{longtable}{@{}p{0.06\textwidth} p{0.33\textwidth} p{0.33\textwidth} p{0.1\textwidth} p{0.08\textwidth}@{}}\toprule Item & Resolução & Gabarito & Veredito & Certo \\ \midrule
a & $\eta$\_A = 60/150 = 0,40 = 0,8$\cdot$$\eta$\_Carnot $\Rightarrow$ $\eta$\_Carnot = 0,5 $\Rightarrow$ T = 450 K & $\eta$1 = 40\%; T = 450 K & igual & ambos \\
b & Q = 150 - 60 = 90 kW & Q rejeitado = 90 kW & igual & ambos \\
c & Não é realista. São necessários 116,7 kW (elétrica) + 33,3 kW (água) + 45 kW (rejeitados a 450 K) = 195 kW > 150 kW, o que viola a 1ª lei. Pela 2ª lei, Q$_{3}$ $\geq$ 20 kW, logo $\dot W$ $\leq$ 85 kW ($\leq$ 255 residências). Aquecendo a água, B fornece 71,7 kW (215 residências), com $\Delta$$\dot S$ $\approx$ +0,044 kW/K, o que é aceitável. & P necessária = 116,7 kW; Q\_água = 33,3 kW; a máquina 2 fornece 71,7 kW < 116,7 kW $\Rightarrow$ NÃO supre as 350 residências (máx. 215); $\Delta$$\dot S$\_universo = +0,04 kW/K $\Rightarrow$ não viola a 2ª lei & igual & ambos \\
\bottomrule\end{longtable}%
\begin{itemize}%
\item%
\textbf{a)} Refiz a conta: $\dot W$ = 80$\cdot$3/4 = 60 kW, $\eta$ = 0,4, $\eta$\_Carnot = 0,4/0,8 = 0,5 e T = 900$\cdot$(1 - 0,5) = 450 K. A resolução e o gabarito dão o mesmo resultado. A transcrição confere com o enunciado original (p. 5).%
\item%
\textbf{b)} Pela 1ª lei, Q = 150 - 60 = 90 kW. Os dois chegam ao mesmo valor. A resolução ainda confere a geração de entropia, 1/30 kW/K > 0, coerente com uma máquina irreversível.%
\item%
\textbf{c)} Conferi com python: P = 350$\cdot$240/720 = 116,67 kW; Q\_água = (50/60)$\cdot$4000$\cdot$10 = 33,33 kW; $\dot W$ = 150 - 45 - 33,33 = 71,67 kW; 71,67$\cdot$3 = 215 residências; $\Delta$$\dot S$ = -1/6 + 1/10 + 33,33/300 = 0,0444 kW/K. A conclusão é a mesma do gabarito: com a água aquecida, B entrega 71,7 kW e não supre as 350 residências (no máximo 215), sem violar a 2ª lei nessa configuração. A resolução vai além e acerta dois pontos extras: o total pedido (195 kW) excede a entrada de 150 kW, e o limite reversível dá $\dot W$ $\leq$ 85 kW (Q$_{3}$ $\geq$ 300$\cdot$(1/6 - 1/10) = 20 kW). O gabarito arredonda $\Delta$$\dot S$ para 0,04 e a resolução dá 0,044; a diferença é só de arredondamento.%
\end{itemize}%
\paragraph{Pontos de dúvida}%
\begin{itemize}%
\item%
\textit{(solution)} No item c, a desigualdade de Clausius está escrita com o sinal trocado. As hipóteses adotam $\Sigma$ Q\_i/T\_i $\leq$ 0, com Q\_i positivo quando o calor entra no ciclo; isso daria +150/900 - 45/450 - Q$_{3}$/300 $\leq$ 0. A resolução, porém, escreve -150/900 + 45/450 + Q$_{3}$/300 $\leq$ 0, que é a variação de entropia dos reservatórios (deveria ser $\geq$ 0). Do jeito que está escrita, a desigualdade levaria a Q$_{3}$ $\leq$ 20 kW, mas a resolução conclui Q$_{3}$ $\geq$ 20 kW, que é o correto. É um erro de digitação que inverte o sentido da desigualdade. Não afeta a resposta final ($\dot W$ $\leq$ 85 kW), mas um corretor rigoroso pode descontar um pouco. A mesma mistura de convenções aparece na verificação do item b, onde se usa -150/900 + 90/450 > 0 como geração de entropia.%
\item%
\textit{(key)} O gabarito dá $\Delta$$\dot S$\_universo = +0,04 kW/K. O valor exato é 2/45 $\approx$ 0,044 kW/K. É só arredondamento, sem efeito na conclusão.%
\item%
\textit{(statement)} O enunciado não diz como interpretar 'realísticos'. O gabarito julga pela potência disponível (71,7 kW < 116,7 kW) depois de atender a água. A resolução mostra também a violação da 1ª lei (195 kW > 150 kW) e o teto da 2ª lei (85 kW, 255 residências). As duas leituras levam à mesma conclusão: a proposta de B não é realista.%
\item%
\textit{(solution)} Há duas hipóteses implícitas, ambas razoáveis e presentes também no gabarito: (i) todo o calor rejeitado a 300 K vai para a água, ou seja, Q$_{3}$ = 33,3 kW; (ii) a demanda residencial é a potência média, 240 kWh/720 h = 1/3 kW por residência, sem considerar picos.%
\end{itemize}%
\section*{Questão 6}%
\href{https://poliedroresolve.sistemapoliedro.com.br/vestibulares/ime/2027/ime-2-fase-dia-2-27-10-2026-2-fase-dia-2/dissertativa-6-fisica-2-fase-dia-2-ime-2027}{Resolução do gabarito}\par\smallskip%
\begin{longtable}{@{}p{0.06\textwidth} p{0.33\textwidth} p{0.33\textwidth} p{0.1\textwidth} p{0.08\textwidth}@{}}\toprule Item & Resolução & Gabarito & Veredito & Certo \\ \midrule
a & v = 68 m/s & v = 68 m/s & igual & ambos \\
b & R = 5$\surd$3 $\approx$ 8,66 m (cos$\theta$ = 1/2) & $\theta$ = 60$^\circ$ (cos$\theta$ = 1/2); R = 5$\surd$3 m & igual & ambos \\
c & T $\approx$ 1067,8 N ($\alpha$ = 60$^\circ$) & $\alpha$ = 60$^\circ$; T $\cong$ 1,07$\cdot$10$^3$ N & equivalente & ambos \\
d & V\_sub $\approx$ 5,51$\cdot$10$^{-}$$^2$ m$^3$ ($\approx$ 55,1 L), com E $\approx$ 551,24 N & E $\approx$ 551 N; V $\cong$ 55$\cdot$10$^{-}$$^3$ m$^3$ & equivalente & ambos \\
\bottomrule\end{longtable}%
\begin{itemize}%
\item%
\textbf{a)} Em A a velocidade (sentido anti{-}horário, ponto mais baixo) aponta para +x, direto para O, porque AO é tangente em A. Então 5/4 = 340/(340 - v), o que dá v = 68 m/s = v\_som/5. Conferido.%
\item%
\textbf{b)} Em B, v = (-v, 0) e o vetor de B até O é (d, -2R), então a fonte se afasta de O. 10/11 = 340/(340 + v$\cdot$cos$\theta$) dá v$\cdot$cos$\theta$ = 34 e cos$\theta$ = 1/2. A resolução usa d/$\surd$(d$^2$ + 4R$^2$) = 1/2 e o gabarito usa tg$\theta$ = 2R/d = $\surd$3. São formas equivalentes e as duas dão R = 5$\surd$3 m $\approx$ 8,660 m.%
\item%
\textbf{c)} sen$\alpha$ = R/L = $\surd$3/2, logo $\alpha$ = 60$^\circ$. T = mv$^2$/(R$\cdot$sen$\alpha$) = (8/5)$\cdot$(34$^2$/$\surd$3) = (8/5)$\cdot$667,4 $\approx$ 1067,8 N. O valor exato é 1067,87 N. A diferença em relação ao gabarito é só de arredondamento.%
\item%
\textbf{d)} E = mg + T$\cdot$cos$\alpha$ = 17,32 + 533,92 $\approx$ 551,2 N. V = E/($\rho$g) $\approx$ 0,0551 m$^3$. É o mesmo valor que 55$\cdot$10$^{-}$$^3$ m$^3$, só escrito em outra notação e com outro arredondamento.%
\end{itemize}%
\paragraph{Pontos de dúvida}%
\begin{itemize}%
\item%
\textit{(statement)} A velocidade de 68 m/s ($\approx$ 245 km/h) para um objeto preso ao fundo de uma piscina é fisicamente irrealista, e por isso a tração resulta enorme (\textasciitilde{}1,07 kN) diante de um peso de \textasciitilde{}17 N. É uma característica dos dados do enunciado e não um erro das respostas, mas vale comentar em sala.%
\item%
\textit{(solution)} As duas soluções usam uma hipótese implícita padrão: o som se propaga no ar com o observador parado, e o efeito Doppler depende só da componente da velocidade da fonte na direção fonte--observador no instante da emissão. A resolução declara isso explicitamente, e o gabarito usa a mesma ideia.%
\item%
\textit{(solution)} Para obter $\alpha$ = 60$^\circ$, a resolução supõe que o raio R (do centro até o objeto, desprezando as dimensões do objeto) é a projeção horizontal do fio de comprimento L preso no fundo, logo abaixo do centro da circunferência. Essa geometria está de acordo com a figura e o gabarito usa a mesma.%
\end{itemize}%
\section*{Questão 7}%
\href{https://poliedroresolve.sistemapoliedro.com.br/vestibulares/ime/2027/ime-2-fase-dia-2-27-10-2026-2-fase-dia-2/dissertativa-7-fisica-2-fase-dia-2-ime-2027}{Resolução do gabarito}\par\smallskip%
\begin{longtable}{@{}p{0.06\textwidth} p{0.33\textwidth} p{0.33\textwidth} p{0.1\textwidth} p{0.08\textwidth}@{}}\toprule Item & Resolução & Gabarito & Veredito & Certo \\ \midrule
a & d = 0,45 m & d = 0,45 m & igual & ambos \\
b & E\_el/|$\Delta$U\_g| = 29/18 $\approx$ 1,6; com sinal, E\_el/$\Delta$U\_g $\approx$ -1,6 & kx $\cong$ 6,44 N; U\_el/$\Delta$U\_g $\cong$ -1,6 & igual & ambos \\
c & v = $\surd$2 $\approx$ 1,4 m/s, quando P2 atinge o vértice C (sentido de B para C) & v2 = $\surd$2 m/s & igual & ambos \\
\bottomrule\end{longtable}%
\begin{itemize}%
\item%
\textbf{a)} Com P2 praticamente fixa, a conservação de energia de P1 dá mg(d0-d)sen30$^\circ$ = Ke$\cdot$q1q2(1/d - 1/d0). Como Ke$\cdot$q1q2 = 0,9 N$\cdot$m$^2$ e mg$\cdot$sen30$^\circ$ = 2 N, a raiz diferente de d = 1 é d = 0,9/2 = 0,45 m. Conferi em python. É de fato um ponto de retorno: em B a repulsão é 0,9 N, menor que 2 N, e em d = 0,45 m vale 4,44 N, maior que 2 N.%
\item%
\textbf{b)} Os dois supõem P2 em equilíbrio na aproximação máxima: kx = 2 + 0,9/0,45$^2$ = 58/9 $\approx$ 6,44 N. Então E\_el/$\Delta$U\_g = ($\frac12$kx$^2$)/(-mg$\cdot$x$\cdot$sen30$^\circ$) = -kx/4 = -29/18 $\approx$ -1,61. A resolução dá o valor -1,6 com sinal, igual ao gabarito, e também o módulo 1,6. O valor de kx (6,44 N) está no meio da conta da resolução (58/9 N).%
\item%
\textbf{c)} Com BC horizontal, a quantidade de movimento ao longo de BC se conserva e as massas são iguais. Assim, P2 percorre 0,15 m até C enquanto P1 recua 0,15 m, e a separação final é 0,75 m. Pela energia, 0,9(1/0,45 - 1/0,75) = 0,8 J = 0,4v$^2$, logo v = $\surd$2 m/s. Conferi numericamente.%
\end{itemize}%
\paragraph{Pontos de dúvida}%
\begin{itemize}%
\item%
\textit{(statement)} Item b: a mola começa relaxada e P2 não é segura. No instante inicial atuam sobre P2 2 N + 0,9 N = 2,9 N sem equilíbrio. Sem amortecimento, uma mola muito rígida oscila rápido em torno do equilíbrio, que muda devagar, com amplitude da ordem de 2,9/k. Nesse caso o valor de x na aproximação máxima não fica determinado: a razão kx/4 iria de cerca de 0,9 a 2,3. O gabarito e a resolução supõem equilíbrio quase estático, o que o 'estime a razão aproximada' do enunciado permite. A resolução declara essa hipótese, e o gabarito a usa sem dizer.%
\item%
\textit{(key)} Convenção de sinal no item b: o gabarito dá -1,6 porque $\Delta$U\_g < 0. Uma resposta só com o módulo (1,6) poderia ser cobrada de outro jeito. A resolução dá as duas formas e coincide com o gabarito.%
\item%
\textit{(key)} Na resolução do gabarito, a frase 'Uma constante que aparece o tempo todo:' vem seguida só de figuras, sem a fórmula. Pelo resto da conta, a constante é mg$\cdot$sen30$^\circ$ = 2 N. É um problema de redação ou de extração e não muda os resultados.%
\item%
\textit{(solution)} No item c, a resolução usa P2C = L0 = 0,15 m e despreza a compressão x $\approx$ 6,44/k. Isso é coerente com 'constante elástica muito alta' e com a hipótese do gabarito.%
\end{itemize}%
\section*{Questão 8}%
\href{https://poliedroresolve.sistemapoliedro.com.br/vestibulares/ime/2027/ime-2-fase-dia-2-27-10-2026-2-fase-dia-2/dissertativa-8-fisica-2-fase-dia-2-ime-2027}{Resolução do gabarito}\par\smallskip%
\begin{longtable}{@{}p{0.06\textwidth} p{0.33\textwidth} p{0.33\textwidth} p{0.1\textwidth} p{0.08\textwidth}@{}}\toprule Item & Resolução & Gabarito & Veredito & Certo \\ \midrule
a & O anteparo está se afastando da rede, porque y\_n = D$\cdot$tan$\theta$\_n cresce com D e as franjas saem pelas bordas. & O anteparo está se afastando da rede. & igual & ambos \\
b & D(t1) = 0,2$\surd$7 $\approx$ 0,53 m & D1 = 0,53 m & igual & ambos \\
c & v = 0,3$\surd$3 - 0,1$\surd$7 $\approx$ 0,254 m/s ($\approx$ 0,25 m/s) & v = 0,254 m/s & igual & ambos \\
d & $\Delta$s = 0,6$\surd$15 - 0,2$\surd$7 $\approx$ 1,79 m & $\Delta$s = 1,792 m & igual & ambos \\
\bottomrule\end{longtable}%
\begin{itemize}%
\item%
\textbf{a)} Mesma conclusão e mesma justificativa: com D maior, os máximos laterais passam de L/2 e o número cai de 7 para 5 e depois para 3.%
\item%
\textbf{b)} d = 2,5 $\mu$m, d/$\lambda$ = 4, sen $\theta$3 = 3/4 e tg $\theta$3 = 3/$\surd$7. De D$\cdot$3/$\surd$7 = 0,6 sai D = 0,2$\surd$7 = 0,529 m (com a tabela, 0,53 m).%
\item%
\textbf{c)} Pela tabela, D2 = 0,6$\surd$3 = 1,038 m e v = (1,038 - 0,530)/2 = 0,254 m/s. Com raízes exatas daria 0,2550 m/s. A diferença vem só de arredondamento, porque o enunciado manda usar a tabela.%
\item%
\textbf{d)} Pela tabela, D1 = 0,6$\surd$15 = 2,322 m e $\Delta$s = 2,322 - 0,530 = 1,792 m. Com raízes exatas daria 1,795 m. A resolução escreve 1,79 m, a mesma coisa com uma casa a menos. O tempo de cerca de 7,1 s após t1 que ela cita é coerente.%
\end{itemize}%
\paragraph{Pontos de dúvida}%
\begin{itemize}%
\item%
\textit{(solution)} A resolução trata a rede como fonte pontual sobre o eixo, desprezando a largura do feixe e da rede. Isso fica implícito no enunciado e no gabarito. A resolução declara essa hipótese, e ela é razoável.%
\item%
\textit{(statement)} No limiar D = L/(2$\cdot$tg$\theta$\_n), a franja $\pm$n fica exatamente na borda do anteparo. Tanto a resolução quanto o gabarito tomam esse instante como o da transição, o que é a leitura natural de 'passa de 7 para 5'.%
\item%
\textit{(key)} Os valores 0,254 m/s e 1,792 m saem das raízes tabeladas. Com raízes exatas dariam 0,255 m/s e 1,795 m. A diferença é só de arredondamento e não muda a nota.%
\end{itemize}%
\section*{Questão 9}%
\href{https://poliedroresolve.sistemapoliedro.com.br/vestibulares/ime/2027/ime-2-fase-dia-2-27-10-2026-2-fase-dia-2/dissertativa-9-fisica-2-fase-dia-2-ime-2027}{Resolução do gabarito}\par\smallskip%
\begin{longtable}{@{}p{0.06\textwidth} p{0.33\textwidth} p{0.33\textwidth} p{0.1\textwidth} p{0.08\textwidth}@{}}\toprule Item & Resolução & Gabarito & Veredito & Certo \\ \midrule
a & u = c($\surd$3 h - L)/($\surd$3 L - h), com impacto em t = ($\surd$3 L - h)/c e R = $\surd$3 h - L & u = c($\surd$3$\cdot$h - L)/($\surd$3$\cdot$L - h) & igual & ambos \\
b & L/$\surd$3 < h < L, ou seja, $\surd$3/3 < h/L < 1 & L/$\surd$3 < h < L & igual & ambos \\
\bottomrule\end{longtable}%
\begin{itemize}%
\item%
\textbf{a)} A expressão é a mesma. Refiz a conta: os versores são d? = ($\surd$3/2, 1/2) e d?' = (0, -1), e a normal externa é n? = (-1/2, -$\surd$3/2), com ângulo de incidência de 30$^\circ$. Escrevendo ct$\cdot$d? = (L, h) + ut$\cdot$n?, sai ut = $\surd$3h - L e ct = $\surd$3L - h. Conferi numericamente com L = 1, h = 0,8: a menor raiz de |ct$\cdot$d? - C| = ut dá t = 0,932/c, u $\approx$ 0,414c, e o raio refletido sai exatamente na direção (0, -1). O sistema de equações é o mesmo do gabarito. A resolução ainda mostra que esse é o primeiro contato, pois g'(t) < 0 na raiz.%
\item%
\textbf{b)} Mesmas condições e mesma conclusão. De ut > 0 sai h > L/$\surd$3; de ct > 0 sai h < $\surd$3L, que é redundante; de u < c, com denominador positivo, sai ($\surd$3+1)h < ($\surd$3+1)L, isto é, h < L. A interseção é L/$\surd$3 < h < L. A resolução acrescenta a leitura geométrica (o centro fica entre a reta a 30$^\circ$ e a reta y = x) e mostra que todo o intervalo é atingível, porque h/L cresce monotonamente com u.%
\end{itemize}%
\paragraph{Pontos de dúvida}%
\begin{itemize}%
\item%
\textit{(statement)} Os extremos ficam de fora (desigualdades estritas). Em h = L/$\surd$3, u = 0 e a esfera degenera num ponto. Em h = L, u = c, que o enunciado exclui (u < c). Uma banca poderia aceitar $\leq$ no extremo inferior se considerasse u = 0 admissível, mas aí o fenômeno não ocorre de fato. Isso não muda a resposta esperada.%
\item%
\textit{(solution)} A resolução declara que a normal externa precisa satisfazer d?$\cdot$n? < 0 (o fóton chega pelo lado de fora). O gabarito usa isso sem dizer. Também despreza o movimento da superfície na reflexão, como o enunciado manda. São hipóteses corretas e estão explícitas.%
\end{itemize}%
\section*{Questão 10}%
\href{https://poliedroresolve.sistemapoliedro.com.br/vestibulares/ime/2027/ime-2-fase-dia-2-27-10-2026-2-fase-dia-2/dissertativa-10-fisica-2-fase-dia-2-ime-2027}{Resolução do gabarito}\par\smallskip%
\begin{longtable}{@{}p{0.06\textwidth} p{0.33\textwidth} p{0.33\textwidth} p{0.1\textwidth} p{0.08\textwidth}@{}}\toprule Item & Resolução & Gabarito & Veredito & Certo \\ \midrule
a & Só a situação 2 (as duas correntes para cima e a partícula descendo). Nas situações 1 e 4, B(0) = 2$\mu$0 i/($\pi$d) $\neq$ 0, então já existe força lateral sobre o eixo. Na situação 3, F\_x = +Kx afasta a partícula do eixo. Na situação 2, F\_x = -Kx, com K = 4$\mu$0 q i v/($\pi$d$^2$), e a força é restauradora. & Somente a situação 2 (1 e 4: B no eixo $\neq$ 0, descartadas; 3: afasta; 2: restaura) & igual & ambos \\
b & p = $\pi$v/$\omega$ $\approx$ ($\pi$d/2)$\cdot$$\surd$($\pi$ m v/($\mu$0 q i)), com $\omega$ = (2/d)$\surd$($\mu$0 q i v/($\pi$ m)) & T = $\pi$d$\cdot$$\surd$(m$\pi$/($\mu$0 i q v)); p = ($\pi$d/2)$\cdot$$\surd$($\pi$mv/($\mu$0 i q)) & igual & ambos \\
\bottomrule\end{longtable}%
\begin{itemize}%
\item%
\textbf{a)} Refiz os sinais. Com fios em x = $\mp$d/2 e as duas correntes para cima, B\_z = ($\mu$0 i/2$\pi$){[}1/(d/2-x) - 1/(d/2+x){]} $\approx$ 4$\mu$0 i x/($\pi$d$^2$). Com F\_x = q v\_y B\_z e v\_y = -v (situação 2), F\_x = -Kx, que é restauradora. Com v\_y = +v (situação 3), a força afasta do eixo. Na simulação RK4 da situação 3, x cresce de forma monotônica a partir de 10$^{-}$$^{7}$. Com correntes opostas (situações 1 e 4), os dois campos entram no papel entre os fios, então B(0) = -2$\mu$0 i/($\pi$d) $\neq$ 0 e a força lateral já age no eixo. As justificativas da resolução coincidem com as do gabarito, inclusive a leitura física: carga descendo equivale a corrente antiparalela, que é repelida pelos fios.%
\item%
\textbf{b)} Partindo de K = 4$\mu$0 q i v/($\pi$d$^2$), obtém{-}se $\omega$$^2$ = K/m e T = 2$\pi$/$\omega$ = $\pi$d$\surd$($\pi$m/($\mu$0 q i v)), o mesmo período do gabarito. Daí p = v$\cdot$T/2 = ($\pi$d/2)$\surd$($\pi$mv/($\mu$0 q i)). Conferi numericamente com RK4 da trajetória exata ($\mu$0 = i = q = m = v = d = 1, desvio inicial pequeno): a distância entre cruzamentos consecutivos do eixo deu 2,78416, e a fórmula dá 2,78416. Ambos tomam p como a distância percorrida em meio período, que é o que mostra a figura (um lóbulo entre cruzamentos consecutivos). A resolução não escreve T explicitamente, mas dá $\omega$, o que é equivalente.%
\end{itemize}%
\paragraph{Pontos de dúvida}%
\begin{itemize}%
\item%
\textit{(statement)} O passo p depende de ler na figura que ele vai de um cruzamento do eixo y ao seguinte, o que corresponde a meio período (p = vT/2). Se alguém o tomasse como o período completo da oscilação, obteria p = vT, o dobro. A figura e o gabarito apoiam T/2, e a resolução justifica essa escolha.%
\item%
\textit{(solution)} A resolução despreza o peso porque g não foi dado. A hipótese é razoável e o gabarito faz o mesmo implicitamente, mas com a velocidade vertical o peso mudaria v ao longo da trajetória. Vale declarar a hipótese, e a resolução declara.%
\item%
\textit{(solution)} Para manter o MHS em x, a resolução aproxima v\_y $\approx$ -v e trata F\_y = -q v\_x B\_z como correção de segunda ordem. A aproximação é correta para amplitude A ? d, e o gabarito usa a mesma hipótese sem comentá{-}la.%
\end{itemize}%
\end{document}